\chapter{Introduction}
The growth of internet has significantly assisted humans, organizations and government in ways no one ever dreamt of ages ago. Internet has 
simplified the lives of individuals and organizations.  Development of robust communication platforms, enhanced economy across the globe, 
elevation in military operations, education, health, entertainment, science, research and invention are some of the privileges provided by 
the internet. Along with the privileges come the burdens in the form of cyberattacks. Intruders tend to overcome the traditional security 
firewalls and benefit by constantly bombarding malicious attacks that the obsolete web security tools and even antiviruses software like 
Kaspersky could never detect \cite{Miani2025}. Thus, urges the need for the deployment of Network Intrusion Detection System (NIDS). Robust AI 
mechanism and machine learning algorithms assist NIDS to detect sophisticated unseen attacks. Thep help in identifying the patterns of the malicious 
attacks analyze them and blocks them enabling a secured network infrastructure in the future. The deployment of Supervised Machine Learning models such 
as Random Forest, XGBoost, Unsupervised Learning models (PCA, K- Means, Autoencoders) and Deep Learning models exhibit better performance in intrusion 
detection from unknown attacks with improved accuracy and F1- Score \cite{deCarvalhoBertoli2023}; \cite{Hamidou2025}; \cite{Wicaksana2026} particularly 
on CICIDS2017 and UNSW-NB15 datasets easily available on public domain. Unlike previous datasets such as KDD Cup 1999 and NSL – KDD, these datasets 
contain real network traffic with varied attacks inhibiting wider traffic patterns \cite{Moustafa2015,Sharafaldin2018}. Even though several AI driven 
approaches have been considered and developed to detect intrusion in the computer network with better accuracy, some models have notable performance on 
the benchmark datasets (CICIDS2017, CSE-CIC-IDS2018 and UNSW-NB15) alone but fail while deployed to detect unknown incoming traffic and thus underperform 
considering their capabilities in cross-dataset generalization across varied computer network. Intruders take an upper hand in this scenario by deploying 
more sophisticated versions of web attacks continuously \cite{Qazi2022}. Some of the reasons why these models fail to generalize conveniently.
\begin{itemize}
    \item Class Imbalance – Most of the web traffics are ‘Genuine’ and few are actual ‘Attack’.
    \item Variable features and size – The training and testing datasets consist of different features; also, the size of the features sourced from different computer network alter.
    \item Computational power – Researchers at an individual level often find it difficult to evaluate the performance of the models given the size of the datasets and the computational power of the device.
\end{itemize}
These issues motivate the researchers to find alternative ways to build robust NIDS to safeguard the computer network from intruders. This allows researchers 
to learn new patterns of attacks and dynamically adjust the models from a classical approach to a revolutionary approach \cite{Kottilingal2024}.
